\section{One actual open for nearby original normal-module freeness}

\begin{theorem}[The same genuine open makes every nearby entire original associated module free]
\label{hir:relative-normal-coefficient-algebra-nearby-associated-free-open}
\source{wlodarczyk05-hironaka}
\lean{Hironaka.HironakaRelativeNormalCoefficientAlgebraNearbyFreeness.exists_original_nearby_associated_free_open}
\leanok
\uses{hir:relative-normal-original-common-domain-coefficient-initial-family,hir:relative-normal-coefficient-uniform-denominator-free-images,hir:relative-normal-graded-internal-finite-relations}
An original coherent ideal $I$ supplies an actual open set $U$ meeting
every neighbourhood of zero, such that at every $a\in U$ the entire
associated graded module of the actual translated original ideal is free over
the actual base holomorphic germ ring.
\end{theorem}
\begin{proof}
Take the internally supplied common domain $D$, genuine base neighbourhood
$W$ and finite polynomial ideal $H$ presenting every nearby entire initial
ideal. The uniform standard-monomial construction gives $b\in D$, $b\ne0$.
Choose a representative $c$ of $b$ in the actual coefficient-function
subalgebra. Its actual function $g$ has nonzero whole origin germ: otherwise
$c$ would belong to the entire origin kernel and its quotient class $b$
would be zero. Put $U=\{a\in W:g(a)\ne0\}$. Holomorphicity on $W$
makes $U$ open. If an origin neighbourhood missed $U$, $g$ would vanish
on its intersection with $W$, contradicting its nonzero whole origin germ.
At every $a\in U$, the actual nearby image of $b$ is the translated
whole germ of $g$, whose value at zero is $g(a)\ne0$; it is therefore
a unit. The uniform free coefficient-image theorem gives the entire nearby
polynomial quotient free over the actual base germ ring. Its whole ideal is
the entire original initial ideal, and the genuine original symbol quotient
equivalence transfers freeness to the entire actual associated module.
\end{proof}

\begin{lemma}[Every original normal layer is finite]
\label{hir:relative-normal-original-normal-layer-finite}
\source{wlodarczyk05-hironaka}
\lean{Hironaka.HironakaRelativeNormalCoefficientAlgebraNearbyFreeness.moduleFinite_originalNormalGradedPiece}
\leanok
\uses{hir:relative-normal-graded-homogeneous-surjection}
For every original germ ideal and every natural $q$, its actual normal
degree-$q$ layer is finite over the base holomorphic germ ring.
\end{lemma}
\begin{proof}
The degree-$q$ homogeneous polynomial submodule is generated by the finitely
many monomials in $t$ variables of total degree $q$. The true homogeneous
symbol map surjects from this finite module onto the entire actual original
normal layer. Its image is consequently finite.
\end{proof}

\begin{theorem}[Every original normal degree is finite free on that same open]
\label{hir:relative-normal-coefficient-algebra-all-degree-free-open}
\source{wlodarczyk05-hironaka}
\lean{Hironaka.HironakaRelativeNormalCoefficientAlgebraNearbyFreeness.all_originalNormalGradedPiece_finite_free_of_associated_free,Hironaka.HironakaRelativeNormalCoefficientAlgebraNearbyFreeness.exists_original_nearby_all_degree_finite_free_open}
\leanok
\uses{hir:relative-normal-coefficient-algebra-nearby-associated-free-open,hir:relative-normal-original-normal-layer-finite}
The same original-input open $U$ makes every actual normal degree-$q$
layer finite free over the actual base germ ring at every $a\in U$, for
all natural $q$ simultaneously.
\end{theorem}
\begin{proof}
Each degree layer is a direct summand of the entire associated graded direct
sum, by its genuine degree inclusion and component projection. A free entire
associated module therefore makes each layer projective, hence flat.
The preceding finite-layer lemma and the local nature of the holomorphic germ
ring make each finite flat layer free. The same $U$ supports the entire free
module, so it supports this argument for every $q$ simultaneously; no
degree-dependent shrinking is used.
\end{proof}

\begin{theorem}[Every original normal truncation is finite free on the same open]
\label{hir:relative-normal-coefficient-algebra-all-truncation-free-open}
\source{wlodarczyk05-hironaka}
\lean{Hironaka.HironakaRelativeNormalCoefficientAlgebraNearbyFreeness.all_originalNormalTruncation_finite_free_of_graded_free,Hironaka.HironakaRelativeNormalCoefficientAlgebraNearbyFreeness.exists_original_nearby_all_truncation_finite_free_open}
\leanok
\uses{hir:relative-normal-coefficient-algebra-all-degree-free-open,hir:relative-normal-uniformtrunc-original-short-exact}
The same original-input open $U$ makes every actual original quotient by the
sum of the original ideal and the $q$th normal-ideal power finite free over
the actual base germ ring, for every point of $U$ and every natural $q$.
\end{theorem}
\begin{proof}
For each nearby original ideal, induct on $q$. The zero truncation is zero,
since the zeroth normal-ideal power is the entire ring. The genuine original
filtration gives a short exact sequence with degree-$q$ layer on the left,
the $(q+1)$st truncation in the middle, and the $q$th truncation on the
right. Both end terms are finite free, by the preceding all-degree theorem
and the inductive hypothesis. The right term is projective, so this short
exact sequence splits; equivalently the free short-exact-sequence theorem
gives a free middle term. The finite-generation exact-sequence theorem gives
a finite middle term. The induction is carried out at every point of the
same $U$ and makes no further shrinking.
\end{proof}
